\chapter{Reset monoids and the flip-flop}\label{ch:reset}

\begin{definition}[Reset monoid]\label{def:resets}
  \lean{KRTheory.TransMon.Resets, KRTheory.TransMon.resetMonoid}
  \uses{def:transmon}
  $U(X) := (X, \{\mathrm{id}\} \cup \{\mathrm{to}_x : x \in X\})$ with
  $a \cdot \mathrm{to}_y = \mathrm{to}_y$ and $a \cdot \mathrm{id} = a$;
  $\mathrm{to}_x$ acts as the constant map $x$.
\end{definition}

\begin{definition}[Flip-flop]\label{def:flipflop}
  \lean{KRTheory.TransMon.flipFlop}\uses{def:resets}
  The flip-flop is $U(\mathbf{2})$: two states, three elements — the
  unique non-group prime of Krohn--Rhodes theory. The monoid part of
  the wreath product of two flip-flops has $3^2 \cdot 3 = 27$ elements.
\end{definition}

\begin{lemma}[Split]\label{lem:reset-split}
  \lean{KRTheory.TransMon.reset_split}
  \uses{def:resets,def:flipflop,def:wreath,def:sdiv}
  For finite $X$, $x_0 \in X$ with $Y := X \setminus \{x_0\}$ nonempty:
  $U(X) \prec U(Y) \wr U(\mathbf 2)$.
\end{lemma}
\begin{proof}
  The flag records ``the current state is $x_0$'':
  $\varphi(y, b) = x_0$ if $b$ else $y$. The covering submonoid consists
  of constant-front pairs $(f, r)$ with two side conditions —
  $r = \mathrm{to}_{\mathrm{false}}$ forces the front value to be a
  reset, and $r = \mathrm{id}$ forces it to be $\mathrm{id}$ — excluding
  exactly the wreath elements with no $U(X)$ counterpart. The covering
  map sends $(g, \mathrm{to}_{\mathrm{true}}) \mapsto \mathrm{to}_{x_0}$
  and $(\mathrm{to}_y, \mathrm{to}_{\mathrm{false}}) \mapsto
  \mathrm{to}_y$.
\end{proof}

\begin{lemma}[DKS 2.12]\label{lem:reset-div-flipflops}
  \lean{KRTheory.TransMon.reset_div_flipFlops}
  \uses{lem:reset-split,lem:wreath-mono,lem:wreathList-append,
        lem:wreath-trivial,def:wreathList}
  For finite nonempty $X$ there is $n$ with
  $U(X) \prec \underbrace{U(\mathbf 2) \wr \cdots \wr
  U(\mathbf 2)}_{n}$.
  Nonemptiness is necessary: an empty-state transformation monoid
  strongly divides only empty-state ones.
\end{lemma}
\begin{proof}
  Strong induction on $|X|$. Base $|X| = 1$: collapse the flip-flop onto
  the point. Step: Lemma~\ref{lem:reset-split}, the induction
  hypothesis, monotonicity, and the append lemma.
\end{proof}
